\section{Related Work}
\label{sec:related_work}

\paragraph{Speculative Decoding Algorithms.}
Speculative decoding accelerates autoregressive generation by decoupling token proposal from verification. Building on early blockwise methods~\citep{stern2018blockwiseparalleldecodingdeep,sun-etal-2021-instantaneous,ge2022lossless,xia2023speculativedecodingexploitingspeculative}, modern approaches employ rejection sampling to exactly preserve the target model's distribution~\citep{chen2023accelerating,pmlr-v202-leviathan23a}. Because inference speedup directly depends on the drafter's efficiency and accuracy, extensive research has focused on optimizing its architecture. Beyond using standalone small language models~\citep{chen2023accelerating,pmlr-v202-leviathan23a}, subsequent work integrates multi-token heads or feature extrapolators directly into the target model~\citep{pmlr-v235-cai24b,ankner2024hydrasequentiallydependentdraftheads,pmlr-v235-li24bt,li-etal-2024-eagle,li2026eagle,zhang2025learningharmonizedrepresentationsspeculative,pmlr-v235-gloeckle24a,liu2024deepseek,cai2025fastmtpacceleratingllminference,eldenk2026attentiondriftautoregressivespeculative}. Other strategies include self-speculation via early exits~\citep{Zhang_2024,Elhoushi_2024,NEURIPS2024_16336d94,xia2025swiftontheflyselfspeculativedecoding}, dynamic vocabulary compression~\citep{zhao2025fr,williams2026speculative}, prompt lookup~\citep{saxena2023prompt,somasundaram2024pldacceleratingllminference}, suffix automata~\citep{hu2024samdecodingspeculativedecoding}, and retrieval~\citep{he2023rest,shen2026draftlessretrievemore}.
To remove the sequential bottleneck of the drafting itself, a line of research proposes parallel or blockwise generation, including Medusa~\citep{pmlr-v235-cai24b}, P-EAGLE~\citep{hui2026peagleparalleldraftingeaglescalable}, PARD~\citep{an2026pard,an2026pard2}, DART~\citep{liu2026dartdiffusioninspiredspeculativedecoding} and DFlash~\citep{chen2026dflash}.
DDTree, TAPS and JetSpec then extend the draft chain to verifiable trees~\citep{ringel2026ddtree,wang2026tapstargetawareprefixtree,hu2026jetspecbreakingscalingceiling}. 
Concurrent efforts include: Domino~\citep{huang2026dominodecouplingcausalmodeling} introduces a CausalEncoder conceptually similar to our RNN Head; DFlare~\citep{zhang2026dflarescalingdraftcapacity} addresses conditioning bottlenecks via layer-wise fusion.

\paragraph{System-Aware Scheduling for Speculative Decoding.}
Beyond drafter architecture, another line of work focuses on determining the optimal number of speculative tokens to generate or verify in each round. To this end, various approaches adapt draft lengths on the fly using confidence heuristics~\citep{mamou2024dynamicspeculationlookaheadaccelerates,li-etal-2024-eagle, du2024glide,liu2026talon,wen2026specboundadaptiveboundedselfspeculation}, learned acceptance predictors~\citep{huang2024specdec++, zacks9172026automtpvllm}, or bandit-style policies~\citep{liu2026notabanditprovablynoregretdrafter}. Furthermore, recognizing speculative decoding as inherently a system-level scheduling problem, recent works optimize overall goodput and latency by adjusting speculation budgets according to real-time system load and request priority~\citep{angelslim2026dcut, liu2024turbospec,10.1145/3772052.3772239,li2026nightjardynamicadaptivespeculative,wu-etal-2025-tetris,hu2026echo,miao2024specinfer,sadhukhan2025magicdecbreakinglatencythroughputtradeoff}.

\paragraph{Parallel Generation.}
Models that generate tokens in parallel offer a decoding latency nearly independent of output length, making them an attractive alternative to autoregressive decoding.
Non-Autoregressive Transformers \citep[NATs,][]{gu2018nonautoregressive} pioneered this direction by predicting all positions independently in a single pass.
However, this forces the model to average over all plausible modes, often producing outputs that mix fragments from different valid sequences. 
Two broad lines of work have emerged to address this limitation.
One direction retains the single-pass architecture but changes what the model sees or how it is trained: introducing latent variables as conditioning input to steer all positions toward a consistent output~\citep{gu2018nonautoregressive,pmlr-v80-kaiser18a,ma-etal-2019-flowseq}, or relaxing the training objective so that the model focuses on producing a single coherent output rather than modeling the full distribution over all valid alternatives~\citep{shao-etal-2021-sequence,pmlr-v139-du21c,qian-etal-2021-glancing,shao2023beyond}.
The other direction reintroduces limited sequential dependency through iterative re-prediction~\citep{ghazvininejad-etal-2019-mask,austin2021structured,li2022diffusionlm}, block-level autoregression~\citep{wang-etal-2018-semi-autoregressive,arriola2025block}, or structured output layers such as CRF~\citep{NEURIPS2019_74563ba2}, CTC~\citep{libovicky-helcl-2018-end,saharia-etal-2020-non}, HMM~\citep{pmlr-v162-huang22m}, and PCFG~\citep{gui2023nonautoregressive}.

Speculative decoding places a further demand that the drafter must provide exact per-token probabilities for the rejection sampling rule. 
Most techniques above cannot readily provide such probabilities due to iterative refinement, latent marginalization, or global normalization.
For instance, in a design closely related to ours, CRF-NAT~\citep{NEURIPS2019_74563ba2} also places a sequential module over parallel hidden states, but its globally normalized partition function prevents exact per-token probability computation. Similarly, when adapting the CTC output layer to parallel speculative decoding, CTC-drafter~\citep{NEURIPS2024_a79054a9} is restricted to greedy verification due to the latent marginalization of alignment paths. DSpark circumvents these limitations by keeping the sequential correction local, so per-token probabilities remain exact softmax evaluations. 



